\chapter*{Glossar}
\addcontentsline{toc}{chapter}{Glossar}

\begin{description}
    \item[ACID] Ein Akronym für Atomicity, Consistency, Isolation, Durability. Diese Eigenschaften garantieren die Zuverlässigkeit von Datenbanktransaktionen.
    \item[Agentic AI] Ein Paradigma der künstlichen Intelligenz, bei dem KI-Systeme (Agenten) eigenständig Ziele verfolgen, Handlungen planen und externe Werkzeuge nutzen.
    \item[API] Application Programming Interface. Eine definierte Schnittstelle zur Kommunikation zwischen verschiedenen Softwarekomponenten.
    \item[BM25] Best Matching 25. Eine Ranking-Funktion, die von Suchmaschinen verwendet wird, um die Relevanz von Dokumenten für eine Suchanfrage basierend auf Begriffshäufigkeiten zu bewerten.
    \item[CSS] Cascading Style Sheets. Eine Stylesheet-Sprache zur Gestaltung von HTML-Dokumenten und Benutzeroberflächen (z.\,B. in GTK).
    \item[Docker] Eine Software zur Containerisierung von Anwendungen, die eine isolierte und reproduzierbare Ausführungsumgebung bietet.
    \item[Embedding] Eine numerische Repräsentation von Informationen (z.\,B. Wörter oder Sätze) in einem hochdimensionalen Vektorraum, die semantische Ähnlichkeiten geometrisch abbildet.
    \item[GTK] GIMP Toolkit. Ein freies Toolkit zur Erstellung von grafischen Benutzeroberflächen, primär für Linux-Umgebungen.
    \item[HNSW] Hierarchical Navigable Small World. Ein effizienter Algorithmus für die approximative Suche nach nächsten Nachbarn (Approximate Nearest Neighbor Search) in Vektorräumen.
    \item[LangChain] Ein Framework zur Entwicklung von Anwendungen, die auf großen Sprachmodellen basieren, insbesondere für die Orchestrierung von Agenten und Tools.
    \item[LLM] Large Language Model. Ein auf riesigen Textmengen trainiertes neuronales Netz, das in der Lage ist, menschenähnliche Sprache zu verstehen und zu generieren.
    \item[MCP] Model Context Protocol. Ein offener Standard zur sicheren und strukturierten Anbindung von KI-Modellen an externe Datenquellen und Werkzeuge.
    \item[Next.js] Ein auf React basierendes Web-Framework, das Funktionen wie serverseitiges Rendering und statische Seitengenerierung bietet.
    \item[NLP] Natural Language Processing. Ein Teilbereich der Informatik und KI, der sich mit der automatischen Verarbeitung natürlicher Sprache beschäftigt.
    \item[Qdrant] Eine spezialisierte Open-Source-Vektordatenbank zur Speicherung und effizienten Suche von Embeddings.
    \item[RAG] Retrieval-Augmented Generation. Ein Verfahren, bei dem Sprachmodelle mit externen Informationen aus einer Datenbank angereichert werden, um präzisere Antworten zu generieren.
    \item[ReAct] Reasoning and Acting. Ein Framework für KI-Agenten, bei dem das Modell explizite Gedankenschritte (Reasoning) mit der Nutzung von Werkzeugen (Acting) kombiniert.
    \item[REST] Representational State Transfer. Ein Architekturstil für zustandslose Netzwerkschnittstellen, die standardisierte HTTP-Methoden verwenden.
    \item[SPECTER2] Ein spezialisiertes Embedding-Modell für wissenschaftliche Publikationen, das auf der Transformer-Architektur basiert.
    \item[Transformer] Eine neuronale Netzwerkarchitektur, die auf dem Attention-Mechanismus basiert und die moderne Sprachverarbeitung (NLP) revolutioniert hat.
\end{description}
